% Shared preamble for the qurl-theory reports.
% Build either report from this folder with:  latexmk -pdf <file>.tex
% (or: pdflatex <file>; bibtex <file>; pdflatex <file>; pdflatex <file>)

\usepackage[T1]{fontenc}
\usepackage{lmodern}
\usepackage[margin=1in]{geometry}
\usepackage{microtype}
\usepackage{amsmath,amssymb,amsthm,mathtools}
\usepackage{booktabs,tabularx,array}
\usepackage{enumitem}
\usepackage[dvipsnames]{xcolor}
\usepackage[authoryear,round]{natbib}
\usepackage{xurl}  % lets long URLs break anywhere (loads url)

% Allow a line break (without an inserted hyphen) after each \_ in text, so long file names can wrap.
\makeatletter
\DeclareRobustCommand{\_}{\ifmmode\nfss@text{\textunderscore}\else\textunderscore\allowbreak\fi}
\makeatother
\usepackage[colorlinks=true,linkcolor=NavyBlue,citecolor=ForestGreen,urlcolor=RoyalBlue]{hyperref}
\usepackage[capitalise,noabbrev]{cleveref}

% ---- theorem-like environments (defined after cleveref so shared counters get the right names)
\theoremstyle{plain}
\newtheorem{theorem}{Theorem}
\newtheorem{proposition}[theorem]{Proposition}
\newtheorem{lemma}{Lemma}
\newtheorem{corollary}{Corollary}
\theoremstyle{definition}
\newtheorem{definition}{Definition}
\newtheorem*{assumptionU}{Assumption U}
\theoremstyle{remark}
\newtheorem{remark}{Remark}

% ---- status tags (robust, so they can appear in section titles and the table of contents)
\DeclareRobustCommand{\statustag}[2]{\textsf{\footnotesize\textcolor{#1}{[#2]}}}
\DeclareRobustCommand{\tagchecked}{\statustag{ForestGreen}{checked}}
\DeclareRobustCommand{\tagproved}{\statustag{RoyalBlue}{proved}}
\DeclareRobustCommand{\tagderived}{\statustag{RoyalBlue}{derived}}
\DeclareRobustCommand{\tagconj}{\statustag{BurntOrange}{conjecture}}
\DeclareRobustCommand{\tagopen}{\statustag{BrickRed}{open}}

% ---- research-problem headings (numbered P1, P2, ... by hand, listed in the table of contents)
\newcommand{\researchproblem}[2]{%
  \subsubsection*{#1.\ #2}\addcontentsline{toc}{subsubsection}{#1.\ #2}}

% ---- math and text macros
\newcommand{\E}{\mathbb{E}}
\newcommand{\R}{\mathbb{R}}
\newcommand{\N}{\mathbb{N}}
\newcommand{\Z}{\mathbb{Z}}
\newcommand{\Unif}{\mathcal{U}}
\DeclareMathOperator{\Var}{Var}
\DeclareMathOperator{\Corr}{Corr}
\DeclareMathOperator{\clip}{clip}
\DeclareMathOperator{\diag}{diag}
\DeclareMathOperator{\sgn}{sign}
\DeclareMathOperator{\fracpart}{frac}
\DeclareMathOperator{\rms}{rms}
\newcommand{\KL}{\mathrm{KL}}
\newcommand{\TV}{\mathrm{TV}}
\newcommand{\script}[1]{\texttt{sanity\_checks/\allowbreak#1}}
\newcolumntype{L}{>{\raggedright\arraybackslash}X}
\setlist{itemsep=2pt,topsep=4pt}

% ---- layout: avoid overfull lines from unbreakable math/citations; ragged-right bibliography
\setlength{\emergencystretch}{1.5em}
\renewcommand{\bibfont}{\raggedright}
